\SourceLeaf{75}
\Needspace{10\baselineskip}\begin{center}
\textbf{BEMERKUNGEN ZU DEN BEIDEN VORSTEHENDEN FRAGMENTEN ÜBER EINEN EULERSCHEN SATZ.}
\end{center}

\noindent\hspace*{\parindent}Das hier von \textsc{Gauss} behandelte \textsc{Euler}sche Theorem besagt folgendes: »Die Function
$V(x, y, p, q, r, s, .\,.\,.)$ von $x, y, p=\dfrac{dy}{dx}$, $q=\dfrac{dp}{dx}$, $r=\dfrac{dq}{dx}$, $s=\dfrac{dr}{dx}$, $.\,.\,.$
liefert stets und nur unter der Bedingung:
\[
N-\dfrac{dP}{dx}+\dfrac{d^2Q}{dx^2}-\dfrac{d^3R}{dx^3}+\dfrac{d^4S}{dx^4}-.\,.\,.=0
\]
ein exactes Differential $Vdx$, wenn hierbei unter $N, P, Q, R, S, .\,.\,.$ die partiellen Ableitungen
\[
\begin{aligned}N=\dfrac{\partial V}{\partial y},\quad P=\dfrac{\partial V}{\partial p},\quad Q=\dfrac{\partial V}{\partial q},\\ R=\dfrac{\partial V}{\partial r},\quad S=\dfrac{\partial V}{\partial s},\ .\,.\,.
\end{aligned}
\]
von $V$ verstanden werden«. \textsc{Euler} hat diesen Satz sowie auch seine Umkehrung im dritten Theile
seiner »Institutiones calculi integralis« und zwar in \S\,92 des die Variationsrechnung behandelnden
Anhanges ausgesprochen und aus den vorangehenden Entwickelungen der Variationsrechnung bewiesen.
\textsc{Gauss} hat seinen vorstehend reproducirten directen Beweis des genannten Satzes in sein
Handexemplar von \textsc{Euler}s Institutiones calculi integralis an der betreffenden Stelle
aufgezeichnet. Den im zweiten Fragmente entwickelten Beweis der Umkehrung hat \textsc{Gauss} in sein
Exemplar von \textsc{Klügels} Lexicon Bd.\,I eingeschrieben.

\begin{flushright}
\textsc{Fricke}.
\end{flushright}

